
\section{Schéma Monolithique}

\begin{frame}{Schéma Monolithique}

\begin{block}{Flux hybrides :$\widetilde{F}$}
On change les flux reconstruit par de nouveaux flux : $\widetilde{F}$
\begin{equation*}
\begin{aligned}
\widetilde{F}_{m+\frac{1}{2}} &=  F^{FV}_{m+\frac{1}{2}} + \theta_{m+\frac{1}{2}} (\widehat{F}_{m+\frac{1}{2}} -F^{FV}_{m+\frac{1}{2}})\\
&=  F^{FV}_{m+\frac{1}{2}} + \theta_{m+\frac{1}{2}} \Delta F_{m+\frac{1}{2}} 
\end{aligned}
\implies
\partial_t \overline{u}^i_m = \frac{-1}{|S_m^i|} (\widetilde{F}^i_{m+\frac{1}{2}}-\widetilde{F}^i_{m-\frac{1}{2}})
\end{equation*}
\end{block}

\begin{block}{critère CFL et solutions de Riemann hybrides}
% Si l'on discrétise en temps par la méthode d'Euler Explicite : 
\begin{tabular}{p{8.5cm}c}
\multirow{4}{9cm}{
$\overline{u}^{i,n+1}_m = \overline{u}^i_m -\frac{\Delta t}{|S_m^i|} (\widetilde{F}^i_{m+\frac{1}{2}}-\widetilde{F}^i_{m-\frac{1}{2}})$\\
$\overline{u}^{i,n+1}_m = \left( 1-  \Delta t_n\frac{\gamma_{m+\frac{1}{2}}+\gamma_{m-\frac{1}{2}} }{|S_m^i|}\right) \overline{{u}_m}^{i,n} + \left(\Delta t_n\frac{\gamma_{m+\frac{1}{2}}+\gamma_{m-\frac{1}{2}} }{|S_m^i|} \right) \widetilde{{u}}_{m+\frac{1}{2}}^-$\\
où $\widetilde{{u}}_{m+\frac{1}{2}}^- = {u}_{m+\frac{1}{2}}^{*,FV} - \theta_{m+\frac{1}{2}} \frac{\Delta F_{m+\frac{1}{2}}}{\gamma_{m+\frac{1}{2}}}$\\
et $\Delta t_n\frac{\gamma_{m+\frac{1}{2}}+\gamma_{m-\frac{1}{2}} }{|S_m^i|} \leq 1$ ce qui nous donne : $\Delta t_n \leq \frac{|S_m^i|}{\gamma_{m+\frac{1}{2}}+\gamma_{m-\frac{1}{2}} }$
}
&
\raisebox{-30mm}{\usebox{\SRiemann}}
\end{tabular}   
\end{block}

\end{frame}
% ------------------------------------------------------------------------------------
% ------------------------------------------------------------------------------------

\begin{frame}{Exemple de critères : Principe du maximum global/local}
\begin{columns}
\begin{column}{0.37\textwidth}
\begin{block}{Principe Maximum global}
\begin{equation}
    u(\cdot,0.)\in [\alpha,\beta] \implies u(\cdot,t^n)\in [\alpha,\beta]
\end{equation}
\end{block}
\end{column}

\begin{column}{0.57\textwidth}
\begin{block}{Principe Maximum Local}
\begin{equation}
    \beta_i^n := \max_{j\in\{i-1,i,i+1\}} u_j^n \geq u_i^{n+1} \geq \min_{j\in\{i-1,i,i+1\}} u_j^n =: \alpha_i^n
\end{equation}
\end{block}
\end{column}

\end{columns}


\begin{block}{critère du Principe du Maximum Local}
    \begin{equation}
\theta_{m+\frac{1}{2}} = \min \left(1, \frac{\gamma}{|\Delta F|}  \left\{
    \begin{aligned}
    \min\left( \beta_{m+1}^n-{u}_{m+\frac{1}{2}}^{*,FV},{u}_{m+\frac{1}{2}}^{*,FV} - \alpha_{m}^n \right), &\text{ si }\Delta F >0  \\
    \min\left( \beta_{m}^n-{u}_{m+\frac{1}{2}}^{*,FV},{u}_{m+\frac{1}{2}}^{*,FV} - \alpha_{m+1}^n \right), &\text{ si }\Delta F <0  
    \end{aligned}\right.
    \right)
\end{equation}  
\end{block}

\begin{block}{cas tests de différents critères}
\red{sinus advection} 
\hspace{2cm}
\red{choc burgers} 
\hspace{2cm}
\red{buckley}

\end{block}

\end{frame}

% ------------------------------------------------------------------------------------
% ------------------------------------------------------------------------------------


\begin{frame}{défaut des critères : \red{erreur induite}, relaxation des critères}

\begin{columns}
\begin{column}{0.68 \textwidth}
\begin{block}{relaxation d'extrema régulier}
\begin{tabular}{p{3cm}c}
\multirow{4}{4cm}{
    $v_i(x) = \partial_x \overline{{u}_{i}} +(x-x_i) \partial_{xx} \overline{{u}_i}$\\
    \vspace{0.1cm}
    $v_{i-1/2}^M := \max(\partial_x\overline{{u}_{i-1}},\partial_x\overline{{u}_{i}})$\\
    \vspace{0.1cm}
    $v_{i-1/2}^m := \min(\partial_x\overline{{u}_{i-1}},\partial_x\overline{{u}_{i}})$\\
    \vspace{0.3cm}
    $v_{i-1/2}^m \leq v_i(x_{i-1/2}) \leq v_{i-1/2}^M$\\
    \vspace{0.1cm}
    $v_{i+1/2}^m \leq v_i(x_{i+1/2}) \leq v_{i+1/2}^M$
}
&
\raisebox{-25mm}{\usebox{\detectextrema}}
\end{tabular}   
\end{block}
\end{column}

\begin{column}{0.3\textwidth}
\begin{block}{convergence des critères}
\red{sinus advection ou burgers}
\end{block}

\end{column}
\end{columns}

\begin{block}{comparaison avec/sans relaxation}
\end{block}

\red{sinus advection} 
\hspace{2cm}
\red{choc burgers} 
\hspace{2cm}
\red{buckley}


\end{frame}

% ------------------------------------------------------------------------------------
% ------------------------------------------------------------------------------------

\begin{frame}{Exemple de critère : critère de positivité}
\begin{columns}

\begin{column}{0.55 \textwidth}
\begin{block}{Espace des solutions des gaz parfaits}
\begin{equation}
    \Omega = \left\{u=(\rho,\rho v,\rho e) \, \biggr \rvert \, \rho >0, \quad \varepsilon= \rho e - \rho \frac{|v|^2}{2} >0\right\}
\end{equation}
\end{block}
\end{column}

\begin{column}{0.35 \textwidth}
\begin{block}{critère de positivité de $\widetilde{\rho}$}
$\theta_{m+\frac{1}{2}}^{\rho>0} = \min \left(1, \frac{\gamma}{|\Delta F^{\rho|} } {\rho}_{m+\frac{1}{2}}^{*} \right)$
\end{block}
\end{column}

\end{columns}

\begin{block}{critère de positivité de $\widetilde{\varepsilon}$}
\begin{tabular}{p{6.2cm} p{6cm}}
$\widetilde{\rho}^{\pm} \widetilde{\rho e}^{\pm} - \frac{|\widetilde{\rho v}^{\pm}|^2}{2}  \geq 0 $ 
\hspace{1cm}$\widetilde{\rho}^{\pm} = \rho^* \pm \theta_{m+\frac{1}{2}} \frac{\Delta F^\rho}{\gamma} $

$\rho^* \, \rho e^* - \frac{|\rho v^* |^2}{2} \geq \mp \theta_{m+\frac{1}{2}} \frac{1}{\gamma}   \left[\Delta F^{\rho} \rho e^* + \Delta F^{\rho e} \rho^* - \Delta F^{\rho v} \rho v^* \right]$
\vspace{0.1cm}
\hspace{2cm}$+ \theta_{m+\frac{1}{2}}^2 \frac{1}{2\gamma^2}\left[2\Delta F^{\rho e}\Delta F^{\rho}- |\Delta F^{\rho v}|^2\right]$
\\
\begin{equation}
\theta_{m+\frac{1}{2}} = \min \left(1, \theta_{m+\frac{1}{2}}^{\rho>0}, \frac{M}{|B|+ min(A,0) } \right)
\end{equation}
&
\multirow{2}{6cm}{
$M = \rho^* \, \rho e^* - \frac{|\rho v|^2}{2}>0 $, 
$B= \frac{1}{\gamma}   \left[\Delta F^{\rho} \rho e^* + \Delta F^{\rho e} \rho^* - \Delta F^{\rho v} \rho v^* \right]$
$A = \frac{1}{2\gamma^2}\left[2\Delta F^{\rho e}\Delta F^{\rho}- |\Delta F^{\rho v}|^2\right]$ 
}
\end{tabular}
\end{block}

\end{frame}

% ------------------------------------------------------------------------------------
% ------------------------------------------------------------------------------------

\begin{frame}{Exemple de critère : critère de positivité}

\red{Sod modifé}\hspace{4cm}\red{Choc-accoustique }\\
\vspace{5cm}
\red{Isentropique}\hspace{4cm}\red{Blast }\\

\end{frame}

% ------------------------------------------------------------------------------------
% ------------------------------------------------------------------------------------


\begin{frame}{Exemple de critère : deuxième critère de positivité (Bonus)}



\end{frame}
